% Section 6 - Atomic operations
% Roberto Masocco <roberto.masocco@uniroma2.it>
% April 22, 2026

% ### Atomic operations ###
\section{Atomic operations}
\graphicspath{{figs/section6/}}

% --- Atomic operations ---
\begin{frame}{Atomic operations}
	\justifying
	An \textbg{atomic operation} is one that completes as an \textbg{indivisible step} from the perspective of all other threads — no other thread can observe a partial result.\\
	\medskip
	Modern processors support atomic \textbg{read-modify-write} instructions at the hardware level (\emph{e.g.}, \texttt{lock xadd} on x86, \texttt{ldxr}/\texttt{stxr} on ARM).\\
	C++ exposes them through \texttt{std::atomic<T>}.\\
	\medskip
	Atomics are appropriate when
	\begin{itemize}
		\item the shared state is a \textbg{single variable} of a simple type (\texttt{int}, \texttt{bool}, pointer);
		\item the operation on it is \textbg{self-contained} (increment, compare-and-swap, flag set/clear);
		\item you need a \textbg{memory barrier} enforcing an \textbg{order}:
      \begin{itemize}
        \item define \textbg{who wins a race};
        \item ensure that \textbg{all cache-to-memory writes are visible to everyone} from the next line.
      \end{itemize}
	\end{itemize}
	\begin{block}{}
		\centering
		For anything more complex than a single variable, use a \textbf{mutex}.
	\end{block}
\end{frame}
\begin{frame}{Atomic operations}{\texttt{std::atomic<T>}}
	\justifying
	\texttt{std::atomic<T>} wraps a value of type \texttt{T} and guarantees that all operations on it are atomic.\\
	The most common operations are:
	\begin{itemize}
		\item \texttt{store(value)}: atomically write;
		\item \texttt{load()}: atomically read;
		\item \texttt{fetch\_add(n)}: atomically read, add, and write; returns the old value;
		\item \texttt{compare\_exchange()}: atomically read, compare, and overwrite with a different value if the comparison succeeds;
		\item the usual arithmetic operators (\texttt{++}, \texttt{+=}, etc.) are overloaded and atomic.
	\end{itemize}
\end{frame}
\begin{frame}[fragile]{Atomic operations}{\texttt{std::atomic<T>}}
	\begin{columns}
		\column{.9\textwidth}
		\begin{lstlisting}[style=codebeamer, language=C++, caption=Fixing the counter race with \texttt{std::atomic}.]
std::atomic<int> counter{0};
void increment(int n) {
  for (int i = 0; i < n; ++i)
    counter++; // atomic read-modify-write
}

int main() {
  std::thread t1(increment, 1000000);
  std::thread t2(increment, 1000000);
  t1.join();
  t2.join();
  std::cout << counter << std::endl; // Exactly 2000000
}
    \end{lstlisting}
	\end{columns}
\end{frame}
\begin{frame}[fragile]{Atomic operations}{Atomic flags}
	\justifying
	A very common use of atomics is a \textbg{stop flag}: a boolean that one thread sets to signal others to terminate their loops.
	\begin{columns}
		\column{.9\textwidth}
		\begin{lstlisting}[style=codebeamer, language=C++, caption=Using an atomic flag to signal thread termination.]
std::atomic<bool> running{true};
void worker() {
  while (running.load()) { /* work */ }
}

int main() {
  std::thread t(worker);
  /* ... */
  running.store(false); // Signal the worker to stop
  t.join();
}
    \end{lstlisting}
	\end{columns}
\end{frame}
\begin{frame}{Atomic operations}{Atomics vs mutexes}
	\justifying
	\textbg{Atomics}:
	\begin{itemize}
		\item protect a \textbg{single variable};
		\item no waiting, no context switch — \textbg{lock-free};
		\item faster for simple operations (counters, flags);
		\item acting as \textbg{memory barriers}.
	\end{itemize}
	\bigskip
	\textbg{Mutexes}:
	\begin{itemize}
		\item protect an \textbg{arbitrary critical section} involving multiple variables and operations;
		\item a thread that cannot acquire the lock is put to sleep by the OS — \textbg{blocking};
		\item the only correct choice when \textbg{a single variable is not enough}.
	\end{itemize}
\end{frame}
